\scanpage{83}

Dans le cas actuel, on aura donc
\[
\begin{aligned}
(68)\qquad V_J^L\simeq{}&\text{le proespace de la forme }V_J^{\prime L},\text{ relatif}\\
&\text{à l'espace }X'=\unsure{X_{I'}^*},\text{ stratifié par les }X_i',\\
&X_i'=X_i\cap X_{I'}^*\quad(i\in I'),\text{ où }X_J'\text{ est fermé,}\\
&\text{et }X_L'\text{ est ouvert, }\deleted{\unsure{\text{et où}}}
\end{aligned}
\]
(Car on a
\[
V_J^L=V_J\cap X_L^*
=\left(\bigcup_{j\in J}V_j^{X'}\right)\cap
\overbrace{\left(\bigcup_{l\in L}V_{\unsure{\deleted{j}\inserted{l}|j}}^{X'}\right)}^{\inserted{X_L^*\ \text{car }X_L^*\text{ ouvert dans }X'}}
=\bigcup_{j\in J,l\in L}V_{j,l}^{X'}
\]
Notons que
\[
V_{j,l}^{X'}=
\begin{cases}
\varnothing&\text{si }j,l\text{ non comparables}^*\\
V_j^{X'}&\text{si }j=l,\text{ et donc }j,l\in J\cap L\\
V_{j,l}^{X'}\ \text{\unsure{standart|standard}}&\text{si }j<l\\
V_{l,j}^{X'}\ \text{\unsure{standart|standard}}&\text{si }l<j\text{ donc }l,j\in J\cap L\\
&\text{(donc }l\in J,j\in L\text{ !)}
\end{cases}
\]
\margin{$^*$ sous réserve que les ouverts de $X$ soient normaux, voire $X$ métrisable....}
donc
\[
(69)\qquad V_J^L=\bigcup_{\substack{j\in J\\l\in L\\j\leq l}}V_{j,l}^{X'}\qquad\deleted{\text{\unsure{pour}}}
\]
(où on se rappellera que $V_{j,j}^{X'}=V_j^{X'}$)

\margin{\uline{NB} \unsure{Sans|Sur} cette formule, valable si $J\prec L$, on voit que dans ce cas on a (69) (\unsure{équiv.} d'homotopie)
\diagram{\begin{tikzcd}[column sep=large,row sep=small]
V_J^L\arrow[r,hook,"{\unsure{(h)|(w)}}","\sim"']\arrow[d,equal,"{\mathrm{déf}}"']&\underbrace{V_J\cap V_L}_{\underset{(\mathrm{déf})}{=}V_{J,L}}\\
V_J\cap X_L^*&
\end{tikzcd}}}

Pour écrire $V_J^L$ par \uline{recollement} des $V_{j,l}^{X'}$, il faut expliciter les
\[
V_{j,l}^{X'}\cap V_{j',l'}^{X'}
\]
qui sont vides, sauf si $\{j,j',l,l'\}$ est un drapeau. C'est un drapeau \deleted{\unsure{de}} ayant au plus quatre termes \inserted{(distincts)}, \deleted{qui} \inserted{ceux-ci} sont dans $J\cup L$, et il en est au \deleted{moins deux dans $J$, deux dans $L$, et au moins un} moins un dans $J$, au moins un dans $L$. Ce sont donc des drapeaux de la forme
\mspage{82}
